\documentclass[11pt, letterpaper]{article}
% ============================================================
% preamble.tex — Shared LaTeX preamble for Learning to Autolens
% ============================================================
% Contains: packages, physics notation macros, formatting.
% Include in each module via % ============================================================
% preamble.tex — Shared LaTeX preamble for Learning to Autolens
% ============================================================
% Contains: packages, physics notation macros, formatting.
% Include in each module via % ============================================================
% preamble.tex — Shared LaTeX preamble for Learning to Autolens
% ============================================================
% Contains: packages, physics notation macros, formatting.
% Include in each module via \input{../preamble}
% ============================================================

% --- Packages ---
\usepackage[utf8]{inputenc}
\usepackage[T1]{fontenc}
\usepackage{amsmath, amssymb, amsthm}
\usepackage{physics}       % \vb, \grad, \div, \curl, \dd, etc.
\usepackage{mathtools}     % \coloneqq, \prescript
\usepackage{bm}            % \bm for bold math
\usepackage{graphicx}
\usepackage{float}
\usepackage{booktabs}      % Professional tables
\usepackage{hyperref}
\usepackage{cleveref}      % \cref for smart cross-refs
\usepackage{xcolor}
\usepackage[margin=1in]{geometry}
\usepackage{enumitem}      % Custom list formatting
\usepackage{fancyhdr}      % Headers and footers
\usepackage{tcolorbox}     % Colored boxes for key results
\usepackage{listings}      % Code listings

% --- Color scheme ---
\definecolor{codeblue}{RGB}{31, 119, 180}
\definecolor{codegray}{RGB}{128, 128, 128}
\definecolor{keyresult}{RGB}{39, 123, 69}
\definecolor{exercise}{RGB}{178, 102, 0}

% --- tcolorbox environments ---
\newtcolorbox{keyresult}[1][]{
  colback=keyresult!5, colframe=keyresult!70,
  fonttitle=\bfseries, title={Key Result}, #1
}

\newtcolorbox{pythonbox}[1][]{
  colback=codeblue!5, colframe=codeblue!50,
  fonttitle=\bfseries, title={PyAutoLens Implementation}, #1
}

\newtcolorbox{exercisebox}[1][]{
  colback=exercise!5, colframe=exercise!50,
  fonttitle=\bfseries, title={Exercise}, #1
}

% --- Code listing style ---
\lstset{
  language=Python,
  basicstyle=\ttfamily\small,
  keywordstyle=\color{codeblue}\bfseries,
  commentstyle=\color{codegray}\itshape,
  stringstyle=\color{keyresult},
  breaklines=true,
  frame=single,
  framerule=0.5pt,
  rulecolor=\color{codegray!50},
}

% --- Physics macros ---
% Vectors (bold upright)
\newcommand{\vtheta}{\bm{\theta}}       % Image-plane position
\newcommand{\vbeta}{\bm{\beta}}         % Source-plane position
\newcommand{\valpha}{\bm{\alpha}}       % Deflection angle
\newcommand{\vx}{\bm{x}}               % Generic position

% Lensing quantities
\newcommand{\thetaE}{\theta_{\rm E}}    % Einstein radius
\newcommand{\SigmaCr}{\Sigma_{\rm cr}}  % Critical surface density
\newcommand{\Dd}{D_{\rm d}}             % Angular diameter distance to lens
\newcommand{\Ds}{D_{\rm s}}             % Angular diameter distance to source
\newcommand{\Dds}{D_{\rm ds}}           % Angular diameter distance lens-source

% Statistical quantities
\newcommand{\chisq}{\chi^2}             % Chi-squared
\newcommand{\chisqred}{\chi^2_{\rm red}} % Reduced chi-squared
\newcommand{\Like}{\mathcal{L}}         % Likelihood
\newcommand{\Evid}{\mathcal{Z}}         % Evidence
\newcommand{\Post}{\mathcal{P}}         % Posterior

% Abbreviations for references
\newcommand{\CK}{C\&K}                  % Congdon & Keeton (2018)
\newcommand{\NB}{N\&B}                  % Narayan & Bartelmann (1997)
\newcommand{\Night}{Nightingale+18}     % Nightingale, Dye & Massey (2018)

% --- Headers ---
\pagestyle{fancy}
\fancyhf{}
\fancyhead[L]{\small\textit{Learning to Autolens}}
\fancyhead[R]{\small\thepage}
\renewcommand{\headrulewidth}{0.4pt}

% --- Theorem environments ---
\theoremstyle{definition}
\newtheorem{definition}{Definition}[section]
\newtheorem{example}{Example}[section]

% ============================================================

% --- Packages ---
\usepackage[utf8]{inputenc}
\usepackage[T1]{fontenc}
\usepackage{amsmath, amssymb, amsthm}
\usepackage{physics}       % \vb, \grad, \div, \curl, \dd, etc.
\usepackage{mathtools}     % \coloneqq, \prescript
\usepackage{bm}            % \bm for bold math
\usepackage{graphicx}
\usepackage{float}
\usepackage{booktabs}      % Professional tables
\usepackage{hyperref}
\usepackage{cleveref}      % \cref for smart cross-refs
\usepackage{xcolor}
\usepackage[margin=1in]{geometry}
\usepackage{enumitem}      % Custom list formatting
\usepackage{fancyhdr}      % Headers and footers
\usepackage{tcolorbox}     % Colored boxes for key results
\usepackage{listings}      % Code listings

% --- Color scheme ---
\definecolor{codeblue}{RGB}{31, 119, 180}
\definecolor{codegray}{RGB}{128, 128, 128}
\definecolor{keyresult}{RGB}{39, 123, 69}
\definecolor{exercise}{RGB}{178, 102, 0}

% --- tcolorbox environments ---
\newtcolorbox{keyresult}[1][]{
  colback=keyresult!5, colframe=keyresult!70,
  fonttitle=\bfseries, title={Key Result}, #1
}

\newtcolorbox{pythonbox}[1][]{
  colback=codeblue!5, colframe=codeblue!50,
  fonttitle=\bfseries, title={PyAutoLens Implementation}, #1
}

\newtcolorbox{exercisebox}[1][]{
  colback=exercise!5, colframe=exercise!50,
  fonttitle=\bfseries, title={Exercise}, #1
}

% --- Code listing style ---
\lstset{
  language=Python,
  basicstyle=\ttfamily\small,
  keywordstyle=\color{codeblue}\bfseries,
  commentstyle=\color{codegray}\itshape,
  stringstyle=\color{keyresult},
  breaklines=true,
  frame=single,
  framerule=0.5pt,
  rulecolor=\color{codegray!50},
}

% --- Physics macros ---
% Vectors (bold upright)
\newcommand{\vtheta}{\bm{\theta}}       % Image-plane position
\newcommand{\vbeta}{\bm{\beta}}         % Source-plane position
\newcommand{\valpha}{\bm{\alpha}}       % Deflection angle
\newcommand{\vx}{\bm{x}}               % Generic position

% Lensing quantities
\newcommand{\thetaE}{\theta_{\rm E}}    % Einstein radius
\newcommand{\SigmaCr}{\Sigma_{\rm cr}}  % Critical surface density
\newcommand{\Dd}{D_{\rm d}}             % Angular diameter distance to lens
\newcommand{\Ds}{D_{\rm s}}             % Angular diameter distance to source
\newcommand{\Dds}{D_{\rm ds}}           % Angular diameter distance lens-source

% Statistical quantities
\newcommand{\chisq}{\chi^2}             % Chi-squared
\newcommand{\chisqred}{\chi^2_{\rm red}} % Reduced chi-squared
\newcommand{\Like}{\mathcal{L}}         % Likelihood
\newcommand{\Evid}{\mathcal{Z}}         % Evidence
\newcommand{\Post}{\mathcal{P}}         % Posterior

% Abbreviations for references
\newcommand{\CK}{C\&K}                  % Congdon & Keeton (2018)
\newcommand{\NB}{N\&B}                  % Narayan & Bartelmann (1997)
\newcommand{\Night}{Nightingale+18}     % Nightingale, Dye & Massey (2018)

% --- Headers ---
\pagestyle{fancy}
\fancyhf{}
\fancyhead[L]{\small\textit{Learning to Autolens}}
\fancyhead[R]{\small\thepage}
\renewcommand{\headrulewidth}{0.4pt}

% --- Theorem environments ---
\theoremstyle{definition}
\newtheorem{definition}{Definition}[section]
\newtheorem{example}{Example}[section]

% ============================================================

% --- Packages ---
\usepackage[utf8]{inputenc}
\usepackage[T1]{fontenc}
\usepackage{amsmath, amssymb, amsthm}
\usepackage{physics}       % \vb, \grad, \div, \curl, \dd, etc.
\usepackage{mathtools}     % \coloneqq, \prescript
\usepackage{bm}            % \bm for bold math
\usepackage{graphicx}
\usepackage{float}
\usepackage{booktabs}      % Professional tables
\usepackage{hyperref}
\usepackage{cleveref}      % \cref for smart cross-refs
\usepackage{xcolor}
\usepackage[margin=1in]{geometry}
\usepackage{enumitem}      % Custom list formatting
\usepackage{fancyhdr}      % Headers and footers
\usepackage{tcolorbox}     % Colored boxes for key results
\usepackage{listings}      % Code listings

% --- Color scheme ---
\definecolor{codeblue}{RGB}{31, 119, 180}
\definecolor{codegray}{RGB}{128, 128, 128}
\definecolor{keyresult}{RGB}{39, 123, 69}
\definecolor{exercise}{RGB}{178, 102, 0}

% --- tcolorbox environments ---
\newtcolorbox{keyresult}[1][]{
  colback=keyresult!5, colframe=keyresult!70,
  fonttitle=\bfseries, title={Key Result}, #1
}

\newtcolorbox{pythonbox}[1][]{
  colback=codeblue!5, colframe=codeblue!50,
  fonttitle=\bfseries, title={PyAutoLens Implementation}, #1
}

\newtcolorbox{exercisebox}[1][]{
  colback=exercise!5, colframe=exercise!50,
  fonttitle=\bfseries, title={Exercise}, #1
}

% --- Code listing style ---
\lstset{
  language=Python,
  basicstyle=\ttfamily\small,
  keywordstyle=\color{codeblue}\bfseries,
  commentstyle=\color{codegray}\itshape,
  stringstyle=\color{keyresult},
  breaklines=true,
  frame=single,
  framerule=0.5pt,
  rulecolor=\color{codegray!50},
}

% --- Physics macros ---
% Vectors (bold upright)
\newcommand{\vtheta}{\bm{\theta}}       % Image-plane position
\newcommand{\vbeta}{\bm{\beta}}         % Source-plane position
\newcommand{\valpha}{\bm{\alpha}}       % Deflection angle
\newcommand{\vx}{\bm{x}}               % Generic position

% Lensing quantities
\newcommand{\thetaE}{\theta_{\rm E}}    % Einstein radius
\newcommand{\SigmaCr}{\Sigma_{\rm cr}}  % Critical surface density
\newcommand{\Dd}{D_{\rm d}}             % Angular diameter distance to lens
\newcommand{\Ds}{D_{\rm s}}             % Angular diameter distance to source
\newcommand{\Dds}{D_{\rm ds}}           % Angular diameter distance lens-source

% Statistical quantities
\newcommand{\chisq}{\chi^2}             % Chi-squared
\newcommand{\chisqred}{\chi^2_{\rm red}} % Reduced chi-squared
\newcommand{\Like}{\mathcal{L}}         % Likelihood
\newcommand{\Evid}{\mathcal{Z}}         % Evidence
\newcommand{\Post}{\mathcal{P}}         % Posterior

% Abbreviations for references
\newcommand{\CK}{C\&K}                  % Congdon & Keeton (2018)
\newcommand{\NB}{N\&B}                  % Narayan & Bartelmann (1997)
\newcommand{\Night}{Nightingale+18}     % Nightingale, Dye & Massey (2018)

% --- Headers ---
\pagestyle{fancy}
\fancyhf{}
\fancyhead[L]{\small\textit{Learning to Autolens}}
\fancyhead[R]{\small\thepage}
\renewcommand{\headrulewidth}{0.4pt}

% --- Theorem environments ---
\theoremstyle{definition}
\newtheorem{definition}{Definition}[section]
\newtheorem{example}{Example}[section]


\fancyhead[C]{\small\textbf{Module 04: Search Chaining \& the SLaM Pipeline}}

\title{\textbf{Module 04: Search Chaining \& the SLaM Pipeline}\\[0.5em]
\large Theory Companion — Learning to Autolens}
\author{Rodrigo C\'ordova Rosado\\
\small Harvard-Smithsonian Center for Astrophysics\\
\small \texttt{rodrigo.cordova\_rosado@cfa.harvard.edu}}
\date{\today}

\begin{document}
\maketitle

\begin{abstract}
This document covers the theory behind search chaining: prior passing,
the curse of dimensionality in Bayesian lens modeling, and the Source-Light-Mass
(SLaM) pipeline architecture. We also introduce pixelized source
reconstructions and regularization, which are central to the SLaM
approach. References: \Night; Nightingale et al.\ (2021);
Suyu et al.\ (2006).
\end{abstract}

\tableofcontents
\newpage

% ============================================================
\section{The Curse of Dimensionality}
\label{sec:dimensionality}
% ============================================================

For a model with $N$ free parameters, the volume of parameter space
scales as $V \propto L^N$ where $L$ is the typical range per parameter.
The fraction of the prior volume that contains the posterior shrinks
exponentially:
\begin{equation}
\frac{V_{\rm posterior}}{V_{\rm prior}} \sim \left(\frac{\sigma_{\rm post}}{\Delta_{\rm prior}}\right)^N
\label{eq:curse}
\end{equation}

For $N = 25$ parameters with $\sigma/\Delta \sim 0.01$ (1\% of prior range),
the posterior occupies $\sim 10^{-50}$ of the prior volume. Finding this
tiny region without guidance is computationally infeasible.

\subsection{Search Chaining as a Solution}

Search chaining addresses this by decomposing the problem:
\begin{enumerate}[nosep]
  \item \textbf{Search 1}: Fit a simplified model ($N_1 \ll N$ parameters)
  \item \textbf{Prior passing}: Use Search 1 posterior as Search 2 prior
  \item \textbf{Search 2}: Fit the full model, but starting in the correct
    region of parameter space
\end{enumerate}

The effective volume ratio for Search 2 becomes:
\begin{equation}
\frac{V_{\rm posterior}}{V_{\rm informed\,prior}} \gg \frac{V_{\rm posterior}}{V_{\rm uninformed\,prior}}
\label{eq:informed-prior}
\end{equation}

% ============================================================
\section{Prior Passing Mechanics}
\label{sec:prior-passing}
% ============================================================

\subsection{Gaussian Prior from Posterior Samples}

Given posterior samples $\{\eta_i^{(k)}\}_{k=1}^K$ for parameter $\eta_i$
from Search 1, the Gaussian prior for Search 2 is:
\begin{align}
\mu_i &= \text{median}\bigl(\{\eta_i^{(k)}\}\bigr) \label{eq:prior-mean} \\
\sigma_i &= w_i \quad \text{(width modifier from config)} \label{eq:prior-sigma}
\end{align}

The \textbf{width modifier} $w_i$ is set in the PyAutoFit prior
configuration and can be:
\begin{itemize}[nosep]
  \item \textbf{Absolute}: $\sigma_i = w_i$ (e.g., $w = 0.05''$ for centres)
  \item \textbf{Relative}: $\sigma_i = w_i \cdot |\mu_i|$ (e.g., $w = 0.5$ for intensity)
\end{itemize}

\subsection{Instance Fixing}

Setting a component to its maximum-likelihood \textbf{instance} removes
it from the parameter space entirely:
\begin{equation}
N_{\rm free}^{(\text{Search 2})} = N_{\rm total} - N_{\rm fixed}
\end{equation}

This is useful when a component is well-constrained and we want to
focus computational effort on other parameters.

% ============================================================
\section{The SLaM Pipeline}
\label{sec:slam}
% ============================================================

\subsection{Pipeline Architecture}

The SLaM (Source, Light, and Mass) pipeline decomposes lens modeling
into four sequential stages:

\begin{table}[H]
\centering
\caption{SLaM pipeline stages.}
\begin{tabular}{lp{3cm}p{3cm}p{4cm}}
\toprule
Stage & Free & Fixed & Purpose \\
\midrule
\textbf{SOURCE LP} & Lens light, mass (SIE), source (S\'ersic)
  & --- & Initial rough model \\
\textbf{SOURCE PIX} & Source pixelization, mass
  & Lens light & Flexible source reconstruction \\
\textbf{LIGHT LP} & Lens light
  & Mass, source & Refined lens light deblending \\
\textbf{MASS TOTAL} & Mass (PowerLaw)
  & Lens light, source & Final mass model \\
\bottomrule
\end{tabular}
\label{tab:slam-stages}
\end{table}

\subsection{Why Source $\to$ Light $\to$ Mass?}

\begin{enumerate}
  \item \textbf{Source first}: A good source model (especially pixelized)
    is needed to accurately subtract arc emission and isolate the lens light.
  \item \textbf{Light before mass}: Clean lens light subtraction reveals
    the lensed arcs without contamination, enabling precise mass modeling.
  \item \textbf{Mass last}: With source and light locked in, the mass
    search has the simplest likelihood surface and fewest degeneracies.
\end{enumerate}

% ============================================================
\section{Pixelized Source Reconstructions}
\label{sec:pixelization}
% ============================================================

\subsection{The Linear Inversion}

Given a mass model (which defines the ray-tracing $\vtheta \to \vbeta$),
the source-plane brightness can be solved as a \textbf{linear system}.
The observed data $\mathbf{d}$ relates to the source pixel brightnesses
$\mathbf{s}$ via:
\begin{equation}
\mathbf{d} = \mathbf{B}\,\mathbf{F}\,\mathbf{s} + \mathbf{n}
\label{eq:linear-system}
\end{equation}
where $\mathbf{F}$ is the \textbf{lensing operator} (maps source pixels
to image pixels via ray-tracing) and $\mathbf{B}$ is the PSF blurring
operator (\Night\ eq.~8; Suyu et al.\ 2006).

The maximum-likelihood source reconstruction is:
\begin{keyresult}
\begin{equation}
\hat{\mathbf{s}} = \bigl(\mathbf{F}^T \mathbf{C}^{-1} \mathbf{F}
+ \lambda \mathbf{H}\bigr)^{-1} \mathbf{F}^T \mathbf{C}^{-1} \mathbf{d}
\label{eq:source-inversion}
\end{equation}
\end{keyresult}

\noindent where $\mathbf{C}$ is the noise covariance matrix,
$\mathbf{H}$ is the \textbf{regularization matrix}, and $\lambda$
controls the smoothness of the reconstruction.

\subsection{Regularization}

Without regularization ($\lambda = 0$), the source reconstruction would
fit the noise — producing a spiky, unphysical source. Regularization
penalizes pixel-to-pixel brightness variations:
\begin{equation}
\text{Regularization penalty} = \lambda\,\mathbf{s}^T \mathbf{H}\,\mathbf{s}
\label{eq:regularization}
\end{equation}

Common choices for $\mathbf{H}$:
\begin{itemize}[nosep]
  \item \textbf{Constant}: $H_{ij} = -1$ for neighbors, $H_{ii} = N_{\rm neighbors}$
    (penalizes differences between adjacent pixels)
  \item \textbf{AdaptiveBrightness}: $\lambda$ varies with source brightness —
    smoother in faint regions, less smoothing in bright regions
\end{itemize}

\subsection{Mesh Types}

\begin{itemize}[nosep]
  \item \textbf{Rectangular}: uniform grid in source plane — simple but wasteful
  \item \textbf{Delaunay}: triangulation in source plane — adapts to geometry
  \item \textbf{Voronoi}: dual of Delaunay — natural pixel shapes
  \item \textbf{Hilbert image mesh}: places pixels along a space-filling curve
    in the image plane, then maps to source plane — concentrates resolution
    where the source is bright
\end{itemize}

\subsection{The Evidence for Pixelized Sources}

The Bayesian evidence for a pixelized source model includes a term
from the linear inversion (Suyu et al.\ 2006 eq.~19):
\begin{equation}
\ln \Evid = -\frac{1}{2}\chisq + \frac{1}{2}\ln\det(\lambda\mathbf{H})
- \frac{1}{2}\ln\det\bigl(\mathbf{F}^T\mathbf{C}^{-1}\mathbf{F} + \lambda\mathbf{H}\bigr)
+ \text{const}
\label{eq:pix-evidence}
\end{equation}

This automatically balances fit quality ($\chisq$) against model
complexity (the regularization terms), implementing Occam's razor.

% ============================================================
\section{Adapt Images}
\label{sec:adapt}
% ============================================================

\textbf{Adapt images} are brightness maps used to guide the pixelization.
After SOURCE LP, the best-fit lensed source image serves as the adapt
image for SOURCE PIX. Pixels are concentrated where the adapt image is
brightest — ensuring high resolution where the source has structure and
low resolution in empty regions.

This is computed by:
\begin{lstlisting}
adapt_image_maker = al.AdaptImageMaker(result=source_lp_result)
\end{lstlisting}

% ============================================================
\section*{References}
% ============================================================

\begin{itemize}[nosep, leftmargin=*]
  \item Nightingale, J.W., Dye, S.\ \& Massey, R.J.\ (2018), MNRAS, 478, 4738
  \item Nightingale, J.W.\ et al.\ (2021), JOSS, 6, 2825
  \item Suyu, S.H.\ et al.\ (2006), MNRAS, 371, 983
  \item Vegetti, S.\ \& Koopmans, L.V.E.\ (2009), MNRAS, 392, 945
\end{itemize}

\end{document}
